\chapter{Construire une première solution réalisable}
\section{Pourquoi commencer par une référence simple}
Une construction rapide fournit à la fois une solution immédiatement utilisable
et un point de départ commun aux recherches. Elle aide à mesurer ce que
l'amélioration apporte réellement. Le terme \code{greedy} de l'interface désigne
ici une chaîne de règles déterministes : sélection géométrique par grille,
affectation au plus proche, construction du cycle puis amélioration 2-opt.
Ce n'est pas une optimisation exacte de chacun de ces sous-problèmes.

La matrice des distances est calculée une seule fois par exécution et partagée
par les étapes suivantes. Tous les coûts utilisent les distances non arrondies.
Les règles d'égalité sur les indices rendent la construction reproductible sans
avoir à répéter artificiellement la même méthode sur plusieurs graines.

\section{Choisir les stations avec une grille}
On calcule le rectangle englobant le nuage de points, puis
$q=\lceil\sqrt{p}\rceil$. Le rectangle est partagé en $q\times q$ cellules.
Dans chaque cellule non vide, on choisit le point le plus proche du centre de
la cellule. La frontière supérieure de la dernière ligne et de la dernière
colonne est incluse pour ne pas perdre les points extrêmes.

Cette règle couvre spatialement le nuage, mais ne fournit pas nécessairement
$p$ stations : des cellules sont vides, $q^2$ peut dépasser $p$ et l'ajout de
la station imposée modifie le compte. Le programme ajuste donc la sélection.
S'il manque des stations, il ajoute le point dont la distance à la station déjà
sélectionnée la plus proche est maximale. Cette complétion privilégie les zones
encore éloignées d'une station. Si la sélection est trop grande, il retire des
stations jugées redondantes par leur proximité aux autres, en préservant la
station imposée.

Le critère de complétion n'est pas une minimisation exacte de la somme des
accès : c'est une règle de couverture. De même, la grille ignore le coût futur
de l'anneau et le poids $\alpha$. Cette simplification explique pourquoi une
recherche sur les stations peut ensuite réduire sensiblement le coût global.

\section{Affecter les points et ordonner les stations}
Une fois $S$ fixé, chaque point non-station est affecté à la station la plus
proche. Cette étape optimise exactement les affectations conditionnellement à
$S$. Les stations s'affectent à elles-mêmes, même si deux points distincts ont
des coordonnées identiques.

Le cycle démarre au sommet imposé. À chaque étape, la station non visitée la
plus proche de la station courante est ajoutée à la séquence. La dernière
station est ensuite reliée implicitement à la première. La règle du plus proche
voisin évite une énumération de tous les ordres possibles, mais elle peut
conserver des détours liés aux choix effectués plus tôt.

\section{Améliorer le cycle avec 2-opt}
Le mouvement 2-opt enlève deux arêtes $(a,b)$ et $(c,d)$, puis les remplace
par $(a,c)$ et $(b,d)$ en inversant l'ordre du segment intermédiaire.
Pour des distances symétriques, la variation de longueur est
\begin{equation}
\Delta L=d_{ac}+d_{bd}-d_{ab}-d_{cd}.
\end{equation}
Il suffit donc de comparer quatre distances, sans réévaluer tout le cycle
pour chaque paire. L'inversion est retenue si elle raccourcit strictement
le parcours au-delà d'une petite tolérance relative. La recherche reprend
alors sur le nouveau cycle jusqu'à ne plus trouver d'amélioration 2-opt.

La première station reste fixe dans la représentation. Les échanges impliquant
l'arête de fermeture sont autorisés : ignorer cette arête laisserait de côté
des améliorations valides. Un test de non-régression sur un carré vérifie
notamment ce cas.

À stations fixées et pour $\alpha>0$, raccourcir l'anneau réduit le coût total
sans toucher aux affectations. Pour $\alpha=0$, 2-opt améliore seulement la
présentation du cycle ; la valeur de l'objectif est inchangée. La fin de 2-opt
fournit une stabilité vis-à-vis de ce voisinage sur l'ordre des stations, avec
les tolérances numériques du programme. Elle ne prouve pas que le cycle est
un optimum du voyageur de commerce.

\section{Algorithme de construction et coût de calcul}
La procédure peut être résumée sans détailler le code :
\begin{enumerate}
\item Calculer les distances et sélectionner les candidats de la grille.
\item Ajouter la station imposée et ajuster l'ensemble pour obtenir $|S|=p$.
\item Calculer les affectations les plus proches.
\item Construire un cycle par plus proche voisin depuis la racine.
\item Appliquer 2-opt au cycle et vérifier la solution complète.
\end{enumerate}
La matrice demande un temps et une mémoire d'ordre $n^2$. Une affectation
directe demande $O(np)$ comparaisons et le plus proche voisin $O(p^2)$.
Une passe de candidats 2-opt examine $O(p^2)$ échanges ; le nombre de passes
dépend des améliorations réalisées. La version de la grille parcourt les
points pour chaque cellule. Ces choix restent simples à contrôler et adaptés
aux tailles étudiées, sans constituer une démonstration de passage à très
grande échelle.
\bilan{La construction donne une référence déterministe, réalise exactement
les affectations à stations fixées et améliore approximativement l'ordre du
cycle. La sélection reste le principal levier des recherches suivantes.}
